% -*- mode : latex; coding: utf-8 -*-
\chapter{Related Work} \label{sec:related-work}

In this chapter, I discuss currently proposed RH mitigations in detail. Research can be classified into three categories: 
mitigation using Near-Row-Refresh, mitigations via swapping, and other alternatives

\noindent\textbf{Mitigation using NRR.} 
The majority of prior proposals \cite{FlippingBits}, \cite{CBT}, \cite{Hydra}, \cite{TWICE}, 
\cite{Graphene} use counter or probabilistic approach to invoke refresh-based mitigation, where adjacent rows of aggressor row
restore leaked charges by row refreshing. Although refreshing victims is a popular choice, the tracker residing in MC
is unaware of refreshed victims, leading to security breakage such as Half-Double \cite{Half_Double}.  

\noindent\textbf{Mitigation via Swapping Rows.}
In order to overcome the security vulnerability of NRR, Saileshwar et al \cite{RRS} proposed RRS,
a novel mechanism of swapping the aggressor row to another random row in the DRAM bank. 
This breaks the spatial relation between victim and aggressor row, effectively compensating for NRR's weakness. 
However, since an attacker can still access swapped rows, RRS must use a threshold of $Flip_{th}/12$ to ensure required level of security. 
Recently proposed AQUA \cite{AQUA} targets such performance overhead of RRS by setting predefined swap-region inside DRAM. 
Since all aggressor rows are swapped to swap-region where the attacker can't access, AQUA does not suffer from a low threshold
and achieve higher performance. 

\noindent\textbf{Other Alternatives.}
BlockHammer \cite{BlockHammer} is an RH mitigation mechanism by throttling aggressor row from accessing DRAM within tREFW
after reaching a predefined threshold. It uses a bloom filter to classify rows between malign and benign accesses. Although
 it is lightweight and incurs small performance overhead, it is susceptible to a denial-of-service attack where an attacker
 repetitively accesses a shared region in a DRAM such as a shared library or application. This can lead to a serious slowdown of
 the entire system. Another promising area of research is manipulating DRAM architecture for stronger security and higher performance,
 such as CROW \cite{CROW}. CROW adds additional \textit{copy-row} to each subarray. Copy-rows are used to save aggressor rows
 at RH attack as well as speed up read/write by faster voltage sensing. Since rows can only be copied to copy-rows at
 the same subarray, CROW cannot provide a security guarantee for concentrated RH attacks targeting a single subarray. 
